\section{Cơ sở lý thuyết}

\subsection{Mô hình ngôn ngữ lớn và hiện tượng ảo giác}

LLM được huấn luyện để dự đoán token kế tiếp sao cho văn bản sinh ra có xác suất cao. Mục tiêu
này không ràng buộc tính đúng đắn: khi tham số của mô hình không chứa thông tin cần thiết, mô
hình vẫn sinh ra câu trả lời trôi chảy nhưng sai — hiện tượng được gọi là \emph{hallucination}
(ảo giác). Với văn bản quy phạm, rủi ro còn cao hơn vì mô hình đã gặp rất nhiều quy chế của các
trường khác trong dữ liệu huấn luyện, nên nó có xu hướng “trả lời theo mẫu chung” thay vì theo
đúng quy định của một trường cụ thể.

\subsection{Retrieval-Augmented Generation}

RAG \cite{lewis2020} tách bài toán làm hai: một bộ \emph{truy xuất} tìm các đoạn văn bản liên
quan từ kho tri thức, rồi một bộ \emph{sinh} viết câu trả lời \textbf{chỉ dựa trên} các đoạn
đó. Tri thức nằm ngoài tham số mô hình nên có thể cập nhật bằng cách nạp lại dữ liệu, không cần
huấn luyện lại; và vì câu trả lời gắn với đoạn trích cụ thể nên có thể kiểm chứng.

Điểm mấu chốt thường bị bỏ qua: \textbf{chất lượng của RAG bị chặn trên bởi chất lượng truy
xuất}. Nếu đoạn văn chứa đáp án không được lấy về, không mô hình nào cứu được câu trả lời. Vì
vậy phần lớn công sức của đồ án này nằm ở bộ truy xuất chứ không ở prompt.

\subsection{Biểu diễn dày (dense embedding)}

Mô hình embedding ánh xạ một đoạn văn thành vector số thực trong không gian nhiều chiều, sao
cho hai đoạn gần nghĩa thì gần nhau theo độ đo cosine. Ưu điểm là bắt được cách diễn đạt khác
nhau: câu hỏi “học bao lâu thì bị buộc nghỉ” vẫn gần với đoạn viết “hết thời gian đào tạo tối
đa”. Nhược điểm là mờ với chi tiết chính xác — con số, mã học phần, số hiệu điều khoản đều bị
“làm mịn” trong vector.

Hệ thống dùng \code{bge-m3} \cite{bgem3}, mô hình đa ngôn ngữ dựa trên XLM-RoBERTa
\cite{xlmr}, vector 1024 chiều, hỗ trợ tiếng Việt tốt và chạy được trên CPU.

\subsection{Biểu diễn thưa và BM25}

BM25 \cite{bm25} là hàm xếp hạng dựa trên tần suất từ khoá, vẫn là chuẩn mực của tìm kiếm văn
bản sau hơn ba thập kỷ. Điểm của một tài liệu $D$ với truy vấn $Q$:

\[
\mathrm{score}(D,Q) \;=\; \sum_{q \in Q} \mathrm{IDF}(q)\cdot
\frac{f(q,D)\,(k_1+1)}{f(q,D) + k_1\left(1-b+b\,\dfrac{|D|}{\mathrm{avgdl}}\right)}
\]

trong đó $f(q,D)$ là tần suất từ $q$ trong $D$, $|D|$ là độ dài tài liệu, $\mathrm{avgdl}$ là
độ dài trung bình của kho, còn $k_1$ và $b$ là tham số điều chỉnh (thường 1,2 và 0,75). BM25
mạnh ở đúng chỗ dense yếu: nó khớp chính xác mặt chữ, và thành phần IDF tự động cho trọng số
cao với từ hiếm — chính là các thuật ngữ chuyên ngành như “phúc tra”, “nợ đọng”, “hậu kiểm”.

\subsection{Tìm kiếm lai và hợp nhất thứ hạng}

Vì dense và sparse mạnh ở hai mặt bù trừ nhau, thực hành phổ biến là chạy cả hai rồi hợp nhất.
Phương pháp thông dụng nhất là \emph{Reciprocal Rank Fusion} (RRF) \cite{rrf}:

\[
\mathrm{RRF}(d) \;=\; \sum_{i} \frac{w_i}{k + \mathrm{rank}_i(d)}, \qquad k = 60
\]

RRF chỉ dùng \textbf{thứ hạng}, không dùng điểm số, nên không cần chuẩn hoá hai thang điểm hoàn
toàn khác nhau (cosine $\in [-1,1]$ và BM25 $\in [0,+\infty)$). Đây vừa là ưu điểm vừa là nhược
điểm: thông tin về \emph{mức độ} liên quan bị vứt bỏ. Mục~\ref{sec:fusion} trình bày cách hệ
thống giữ lại thông tin đó.

\subsection{Rerank}

Bộ truy xuất thường lấy về nhiều ứng viên hơn mức cần (ví dụ 20) rồi dùng một mô hình mạnh hơn
chấm lại để chọn ra 5 đoạn tốt nhất. Chuẩn mực là \emph{cross-encoder}: đưa cặp (câu hỏi, đoạn
văn) vào cùng một mô hình Transformer để mô hình nhìn thấy cả hai cùng lúc — chính xác hơn hẳn
so với so sánh hai vector độc lập, nhưng chi phí tính toán tỉ lệ với số ứng viên.

Trong bối cảnh đồ án, máy đã phải nuôi một LLM và một mô hình embedding trên CPU; thêm
\code{bge-reranker-v2-m3} nghĩa là thêm PyTorch và khoảng 2\,GB bộ nhớ. Mục~\ref{sec:rerank}
trình bày phương án thay thế không cần mô hình.

\subsection{Thách thức riêng của tiếng Việt và văn bản quy phạm}

\begin{itemize}
  \item \textbf{Ranh giới từ.} Tiếng Việt viết rời theo âm tiết; “học phần” là một từ nhưng là
  hai token khi tách theo khoảng trắng. Các thư viện BM25 có sẵn (ví dụ mô hình
  \code{Qdrant/bm25} của fastembed) áp dụng stemming và danh sách stopword tiếng Anh, làm hỏng
  hoàn toàn tiếng Việt.
  \item \textbf{Dấu thanh.} Người dùng thường gõ không dấu (“hoc phan khong dat”). Chuỗi không
  dấu không khớp bất kỳ token có dấu nào trong chỉ mục.
  \item \textbf{Khoảng cách ngôn ngữ.} Sinh viên dùng tiếng lóng và viết tắt; văn bản dùng
  thuật ngữ pháp lý đầy đủ. “Rớt môn” không xuất hiện một lần nào trong 57 trang quy định.
  \item \textbf{Cấu trúc phân cấp.} Ý nghĩa của khoản 5.2 phụ thuộc vào Điều 5 và Chương 2 chứa
  nó. Cắt văn bản theo cửa sổ ký tự cố định sẽ làm mất ngữ cảnh phân cấp này.
  \item \textbf{Bảng biểu mang dữ liệu chính.} Nhiều đáp án quan trọng nằm trong bảng (thời
  gian đào tạo theo loại hình, thang quy đổi điểm chữ, chuẩn ngoại ngữ). Bảng bị cắt đôi là đáp
  án bị mất.
  \item \textbf{Bản scan.} PDF gốc 19\,MB là ảnh quét, không có lớp văn bản; buộc phải OCR, và
  chú thích nguồn gốc ở lề phải dễ bị trộn lẫn vào thân bài.
\end{itemize}

\chapterbreak
